\input{../tex-inputs/latex-leseansicht-vorspann}

\section[Carl von Torresani an Arthur Schnitzler, 9. 12. 1892]{L04326 Carl von Torresani an Arthur Schnitzler, 9. 12. 1892}
\nopagebreak\mylabel{L04326v}
\rehead{ }\normalsize\beginnumbering\briefempfaengerindex{Schnitzler, Arthur@\textsc{Schnitzler, Arthur}!zzzTorresani-Lanzenfeld, Carl von@\emph{von Carl von Torresani-Lanzenfeld}!1892-12-091@{9. 12. 1892}|(be}
\toendnotes[C]{\smallbreak\pagebreak[2]}
\correspDesc{Versand  durch Carl von Torresani am 9. 12. 1892 in Rom
\newline{}Erhalt  durch Arthur Schnitzler im Zeitraum [10. 12. 1892 – 14. 12. 1892?] in Wien}\toendnotes[C]{\smallbreak}
\Standort{CUL, Schnitzler, B 106.}
\physDesc{Brief, 1 Blatt, 4 Seiten, 2237 Zeichen
\newline{}Handschrift: schwarze Tinte, deutsche Kurrent}\toendnotes[C]{\smallbreak}
\pstart
           \raggedleft{}{\pb}\textsc{Rom Piazza Montecitorio 121\oindex{Piazza Montecitorio 121@\textbf{Piazza Montecitorio 121}, \emph{Wohngebäude}|pw}}\pend
           
\pstart
           \raggedleft{}9. 12. 92.\pend
           
\pstart{}Mein hochverehrter Herr Doktor!\pend\vspace{0.5em}
\pstart
           Auf das allerfreudigſte wurde ich vor nun \textsc{cca} 14 Tagen
               durch das Eintreffen ihres Buches: \textsc{Anatol\pwindex{Schnitzler, Arthur 15. 5. 1862 Wien – 21. 10. 1931 ebd.@\textsc{Schnitzler, Arthur} (15. 5. 1862 Wien – 21. 10. 1931 ebd.), \emph{Schriftsteller, Mediziner}!Anatol@\strich\emph{Anatol}|pw}} überraſcht, für deſſen Sendung und liebenswürdige Widmung ich hiemit meinen
               wärmſten Dank ausſpreche.\pend
           
\pstart
           Ich wollte nicht{ }ſchreiben, bevor ich das Buch\pwindex{Schnitzler, Arthur 15. 5. 1862 Wien – 21. 10. 1931 ebd.@\textsc{Schnitzler, Arthur} (15. 5. 1862 Wien – 21. 10. 1931 ebd.), \emph{Schriftsteller, Mediziner}!Anatol@\strich\emph{Anatol}|pwv} mit der Gründlichkeit u Aufmerksamkeit, wie es als
               Ihre Arbeit verdient, geleſen haben würde. Das iſt nun geſchehen {\pb}und ich kann Ihnen nur ſagen, daß ich
               bedauere, zu Ende zu ſein. Es ist geiſtreich von Anfang zum Schluße. Der{ }ſubtil
               zerſetzende, wühlende, analyſirende Geiſt, der das ganze durchdringt, das Gefühl zum
               Sündenbock der Vernunft macht, erinnert lebhaft an \textsc{Maupaſsant\pwindex{Maupassant, Guy de 5.\,8.\,1850 Tourville-sur-Arques – 7.\,7.\,1893 Paris@\textsc{Maupassant, Guy de} (5.\,8.\,1850 Tourville-sur-Arques – 7.\,7.\,1893 Paris), \emph{Schriftsteller}|pw}} mit dem Sie auch die Kraft und Prägnanz des Ausdrucks gemeinſam haben, –
               während die dialogisierende Form in Ihrer franzöſiſchen\oindex{Frankreich@\textbf{Frankreich}|pw} Feinheit an \textsc{Gyp\pwindex{Mirabeau, Sibylle de 15.\,8.\,1850 Plumergat – 30.\,6.\,1932 Neuilly-sur-Seine@\textsc{Mirabeau, Sibylle de} (15.\,8.\,1850 Plumergat – 30.\,6.\,1932 Neuilly-sur-Seine), \emph{Schriftstellerin}|pw}} anklingt. Als die Perle der Sammlung\pwindex{Schnitzler, Arthur 15. 5. 1862 Wien – 21. 10. 1931 ebd.@\textsc{Schnitzler, Arthur} (15. 5. 1862 Wien – 21. 10. 1931 ebd.), \emph{Schriftsteller, Mediziner}!Anatol@\strich\emph{Anatol}|pwv} möchte ich »Denksteine\pwindex{Schnitzler, Arthur 15. 5. 1862 Wien – 21. 10. 1931 ebd.@\textsc{Schnitzler, Arthur} (15. 5. 1862 Wien – 21. 10. 1931 ebd.), \emph{Schriftsteller, Mediziner}!Denksteine@\strich\emph{Denksteine}|pw}«
               bezeichnen. Der Schluß iſt {\pb}meiſterhaft,
               unerwartet, überwältigend – und ſo pſychologiſch \uline{berechtigt}. — Einige Aphorismen möchte ich wirklich den beſten \strikeout{beste} die ich kenne einreihen;{ }ſo \strikeout{\textcolor{gray}{werden}} zB \introOben{}dieses\introOben{}: »gewiſſe Arten den Wahrheit{ }ſind nur die
               Aufrichtigkeit ermüdeter Lügner«. Der Ausdruck: Hypochonder der Liebe, jener andere:
               »leichtſinniger Melancholiker« entzücken mich durch ihre Richtigkeit. Kurz, die
               Lektüre iſt mir ein Genuß geweſen. Nehmen Sie meine aufrichtigſten Complimente –\pend
           
\pstart
           Was macht Ihr Stück: »Das Märchen\pwindex{Schnitzler, Arthur 15. 5. 1862 Wien – 21. 10. 1931 ebd.@\textsc{Schnitzler, Arthur} (15. 5. 1862 Wien – 21. 10. 1931 ebd.), \emph{Schriftsteller, Mediziner}!Märchen. Schauspiel in drei Aufzügen@\strich\emph{Das Märchen. Schauspiel in drei Aufzügen}|pw}?« Ich habe nie
               verſäumt, in der N.F.P\pwindex{Neue Freie Presse@\emph{Neue Freie Presse}|pw} unter den
               »Theaternachrichten« danach zu{ }ſchauen, habe es aber, bis {\pb}jetzt nicht finden können.\pend
           
\pstart
           Ich{ }ſchicke diesen Brief auf Gut Glück an Ihre alte Adreſſe\oindex{Wien@\textbf{Wien}!I., Innere Stadt@\textbf{I., Innere Stadt}!Kärntnerring 12/Bösendorferstraße 11@\textbf{Kärntnerring 12/Bösendorferstraße 11}, \emph{Wohngebäude}|pwv}, überzeugt daß er Ihnen{ }ſchließlich zukommen wird.
               Bitte mir \label{K_L04326-1v}\edtext{Ihre neue\oindex{Wien@\textbf{Wien}!I., Innere Stadt@\textbf{I., Innere Stadt}!Wohnung und Ordination Arthur Schnitzler Grillparzerstraße 7/3. Stock@\textbf{Wohnung und Ordination Arthur Schnitzler Grillparzerstraße 7/3. Stock}, \emph{Ordination}|pwv}}{\lemma{\textnormal{\emph{Ihre neue}}}\Cendnote{\textnormal{Schnitzler übersiedelte laut \emph{Tagebuch}\pwindex{Schnitzler, Arthur 15. 5. 1862 Wien – 21. 10. 1931 ebd.@\textsc{Schnitzler, Arthur} (15. 5. 1862 Wien – 21. 10. 1931 ebd.), \emph{Schriftsteller, Mediziner}!Tagebuch@\strich\emph{Tagebuch}|pwk} am 15. 10. 1892 von der Giselastraße 11\oindex{Wien@\textbf{Wien}!I., Innere Stadt@\textbf{I., Innere Stadt}!Kärntnerring 12/Bösendorferstraße 11@\textbf{Kärntnerring 12/Bösendorferstraße 11}, \emph{Wohngebäude}|pwk} in die Grillparzerstraße 7\oindex{Wien@\textbf{Wien}!I., Innere Stadt@\textbf{I., Innere Stadt}!Wohnung und Ordination Arthur Schnitzler Grillparzerstraße 7/3. Stock@\textbf{Wohnung und Ordination Arthur Schnitzler Grillparzerstraße 7/3. Stock}, \emph{Ordination}|pwk}.}}}\label{K_L04326-1} möglichſt umgehend mitzutheilen; ich möchte
               Ihnen mein{ }ſoeben erſchienenes »Oberlicht\pwindex{Torresani-Lanzenfeld, Carl von 19.\,4.\,1846 Mailand – 16.\,4.\,1907 Nago-Torbole@\textsc{Torresani-Lanzenfeld, Carl von} (19.\,4.\,1846 Mailand – 16.\,4.\,1907 Nago-Torbole), \emph{Schriftsteller, Offizier}!Oberlicht. Künstlerroman@\strich\emph{Oberlicht. Künstlerroman}|pw}« auf{ }ſichere Weiſe zukommen laſſen. Auch würde ich Ihnen, beſter Herr Doctor,{ }ſehr für die
               Angabe von \textsc{H. Bahrs\pwindex{Bahr, Hermann 19.\,7.\,1863 Linz – 15.\,1.\,1934 München@\textsc{Bahr, Hermann} (19.\,7.\,1863 Linz – 15.\,1.\,1934 München), \emph{Schriftsteller, Kritiker}|pw}} Adreſſe verbunden{ }ſein{[}.{]}\pend
           
\pstart
           Wir{ }ſitzen hier in Wärme u Sonnenſchein, gehen ohne Mäntel aus, u leſen mit
               angenehmem Gruseln Berichte über Froſt u Schneeverwehungen aus den Binnenländern.
               Bitte grüßen sie mir herzl. \textsc{Salten\pwindex{Salten, Felix 6.\,9.\,1869 Budapest – 8.\,10.\,1945 Zürich@\textsc{Salten, Felix} (6.\,9.\,1869 Budapest – 8.\,10.\,1945 Zürich), \emph{Schriftsteller, Journalist, Chefredakteur}|pw}}, \textsc{Bahr\pwindex{Bahr, Hermann 19.\,7.\,1863 Linz – 15.\,1.\,1934 München@\textsc{Bahr, Hermann} (19.\,7.\,1863 Linz – 15.\,1.\,1934 München), \emph{Schriftsteller, Kritiker}|pw}}, \textsc{Loris\pwindex{Hofmannsthal, Hugo von 1.\,2.\,1874 Wien – 15.\,7.\,1929 Rodaun@\textsc{Hofmannsthal, Hugo von} (1.\,2.\,1874 Wien – 15.\,7.\,1929 Rodaun), \emph{Schriftsteller}|pw}}, \textsc{Beraton\pwindex{Bératon, Ferry 6.\,12.\,1859 Wien – 11.\,2.\,1900 Venedig@\textsc{Bératon, Ferry} (6.\,12.\,1859 Wien – 11.\,2.\,1900 Venedig), \emph{Schriftsteller, Journalist, Maler}|pw}}, \introOben{}\textsc{Beer-Hoffmann\pwindex{Beer-Hofmann, Richard 11.\,7.\,1866 Wien – 26.\,9.\,1945 New York City@\textsc{Beer-Hofmann, Richard} (11.\,7.\,1866 Wien – 26.\,9.\,1945 New York City), \emph{Schriftsteller}|pw}},\introOben{} u die \label{T_L04326-1v}\edtext{andern Herrn}{\lemma{\textnormal{\emph{andern Herrn}}}\Cendnote{\textnormal{ab hier bis zum Schluss entlang des linken
                     Seitenrandes geschrieben}}}\label{T_L04326-1}\label{K_L04326-2v}\edtext{des »allerjüngſten«}{\lemma{\textnormal{\emph{des »allerjüngsten«}}}\Cendnote{\textnormal{Anspielung auf den literarischen Nachwuchs Wiens\oindex{Wien@\textbf{Wien}, \emph{Verwaltungsgebiet}|pwk}, für den sich bald »Jung Wien« als Begrifflichkeit durchsetzte}}}\label{K_L04326-2} u laſſen Sie ſich ſelbſt
            herzl. die Hand drücken von Ihrem in wirklicher Verehrung ergebnen\pend
           \pstart \spacefill\mbox{Torresani}\pend{}\selectlanguage{ngerman}\endnumbering\briefempfaengerindex{Schnitzler, Arthur@\textsc{Schnitzler, Arthur}!zzzTorresani-Lanzenfeld, Carl von@\emph{von Carl von Torresani-Lanzenfeld}!1892-12-091@{9. 12. 1892}|)be}\mylabel{L04326h}
\begin{anhang}
\end{anhang}\newcommand{\dateiname}{L04326}\newcommand{\titel}{Carl von Torresani an Arthur Schnitzler, 9. 12. 1892}\newcommand{\editorInnen}{Herausgegeben von Jahnke, SelmaMüller, Martin Anton}\input{../tex-inputs/latex-leseansicht-abspann}
